\chapter{Logico-Computational Rationality: Changing Information and Acting Together}
\label{sec:p08-rationality-2-6}

Van Benthem, Liu and Smets treat logic as both a practical technology for
reasoning and a theoretical investigation of the structure and limits of
rational agency. Their 2021 chapter follows the passage from inference
and computation to information change, games, social networks and machine
learning. Its claim is not that rationality reduces to deductive validity,
or that every rational process must visibly manipulate logical formulas.
It is that observation, correction, communication and interaction are
also candidates for systematic logical analysis
\cite[pp.~185--194]{VanBenthemLiuSmets2021Logico}.

\section{What the logical perspective does and does not claim}

The distinction between a valid inference and a useful action is already
visible in elementary examples. From \(A\to B\) and \(\neg B\), modus
tollens licenses \(\neg A\). From \(A\to B\) and \(\neg A\), however,
\(\neg B\) does not follow: a valuation with \(A\) false and \(B\)
true satisfies both premises but falsifies the proposed conclusion.
An agent who avoids this error has respected a consequence relation;
that fact alone says nothing about the accuracy of the premises, the
cost of obtaining them, or the wisdom of the agent's goals. The chapter
accordingly distinguishes logic's practical role in organizing reasoning
from its foundational role in exposing the structure of a practice.
The latter need not make the practice itself an explicit proof search.
Formal languages and procedures are also cultural artifacts, interacting
with ordinary communication and computational devices rather than
replacing them \cite[pp.~185--186]{VanBenthemLiuSmets2021Logico}.

The historical route runs from ancient inference and the combinatorial
projects of Llull and Leibniz to programmable computation. Babbage and
Lovelace exemplify the move from an isolated calculation to a mechanism
whose behavior can be programmed. An algorithm needs a representation
of its data as well as instructions: a different organization of a
problem can change what a procedure makes feasible. The chapter uses
this history to connect proof, definition and algorithm, not to identify
them. Its later account of computational architectures also distinguishes
an algorithm's organization from its input--output function. Distributed
systems add a question about who computes and communicates, alongside
what is computed and how \cite[pp.~186--188]{VanBenthemLiuSmets2021Logico}.

Turing's conversational test adds a different kind of question: can
external behavior distinguish the machine from a human interlocutor?
The chapter treats this as a performance criterion whose interpretation
has been contested and whose demands can shift, not as an established
equivalence between successful conversation and every internal form of
understanding. Behavioral evidence is also important when assessing
humans; access to another person's reasoning is ordinarily indirect.
The test's human--machine interaction is itself a hybrid social setting.
Agents with different observational, inferential and communicative
capacities need not be instances of one uniform actor type. Recognizing
that diversity calls for revisiting model assumptions; it does not
follow that unequal computational abilities imply unequal moral rights
\cite[p.~187]{VanBenthemLiuSmets2021Logico}.

Limits require their hypotheses. The first-incompleteness and Rosser
results in \cref{thm:godel-first,thm:rosser} concern suitable
effectively axiomatized arithmetic theories with their stated consistency
conditions, not every sufficiently interesting notation. The effective
halting problem in \cref{thm:halting-undecidable} and first-order
validity in \cref{thm:fol-validity-undecidable} have specified encodings
and problem domains. Undecidability of a class does not imply that every
particular instance is beyond solution. Nor does more expressive syntax
by itself establish a uniform monotone complexity ranking across all
languages, fragments and encodings. The source's productive lesson is
to investigate expressiveness, representation and resources together.
Automata and memory restrictions, tractable fragments and verification
procedures remain useful despite general impossibility results.
Universality is a simulation property of a suitably programmed machine,
not a guarantee that all computations halt or finish within a usable
budget. The absence of an analogous agreed universal rational agent
therefore does not contradict the universality of computation.

P12 makes this distinction operational. Its
\cref{def:p12-encoding-audit,def:p12-automated-result-audit} asks which
encoded task was attempted, what calculus and resources were used, and
whether an output is a checked certificate or an inconclusive search.
A timeout is not a countermodel. Conversely, an abstract logical law is
not already a runnable implementation with an acceptable resource cost.
The chapter's theoretical and practical perspectives are complementary
precisely because neither settles the other's questions automatically.

\section{Information change is an operation, not just another premise}

The source's richer agent can observe, infer, remember and communicate;
evaluate an information source; reconsider a fallible belief; and relate
beliefs to preferences, goals and intentions. Robots with noisy sensors
illustrate why reliable knowledge and defeasible belief cannot simply
be interchangeable. A belief--desire--intention architecture describes
different roles within an agent, but naming those roles does not prove
that the resulting agent succeeds, or select a uniquely correct utility
function \cite[pp.~188--190]{VanBenthemLiuSmets2021Logico}.

\begin{hbdefinition}{def:p08-rat26-epistemic-update}{Finite hard-information model}
Let \(M=(W,(\sim_i)_{i\in I},V)\), where \(W\) is a nonempty
finite set, each \(\sim_i\) is an equivalence relation, and \(V\)
assigns truth sets to atoms. \(K_i\psi\) holds at \(w\) when
\(\psi\) holds at every \(i\)-accessible world. \(C\psi\) holds
when \(\psi\) holds at every world reachable by a finite sequence
of any agents' accessibility steps, including the empty sequence.
For \(P=\{w:M,w\models\phi\}\), the truthful public announcement
\(!\phi\) restricts the domain, relations and valuation to \(P\).
The box \([!\phi]\psi\) is true at \(w\) if \(\phi\) is false
there, or if \(\psi\) holds at \(w\) in the restricted model.
An actually executable announcement at \(w\) requires \(w\in P\).
This is the P8 operation in \cref{def:p08-public-announcement}, not
a general private or misleading event.
\end{hbdefinition}

\begin{hbtheorem}{thm:p08-rat26-hard-update}{proved=annex:p08.rat2.6.hard-update}{Knowledge after a public factual update}
In the preceding semantics,
\[
 [!\phi]K_i\psi\ \leftrightarrow
 (\phi\to K_i[!\phi]\psi).
\]
If \(\phi\) is propositional, with no epistemic or dynamic operators,
then \([!\phi]C\phi\) is valid. For the two-world model with
\(W=\{w_+,w_-\}\), \(V(p)=\{w_+\}\), universal \(\sim_1\)
and identity \(\sim_2\), agent 1 initially knows neither \(p\) nor
\(\neg p\), but knows that agent 2 knows which is true. At \(w_+\),
a truthful public reply \(p\) leaves a singleton model with \(Cp\).
\end{hbtheorem}

This supplies the source's question--answer example with an explicit
information structure. An utterance ``\(p\)?'' alone does not logically
entail the questioner's ignorance: examinations and rhetorical questions
are counterexamples. Inferring ignorance requires a sincere-information-
seeking protocol; inferring the other person's prior knowledge requires
a competence or infallible-source assumption. A truthful answer alone
does not establish that the speaker knew beforehand. The public factual
update does establish common knowledge within this model, including
common awareness of the update. It does not establish that these
protocol assumptions hold in an observed conversation.

The restriction to factual \(\phi\) in the common-knowledge conclusion
matters. An epistemic announcement can change its own truth conditions:
``\(p\) and you do not know \(p\)'' may be true before it is announced
and false afterward. The displayed reduction law is one semantic
calculation. The existing \cref{thm:p08-pal-reduction} and
\hyperref[sheet:hilbert-pal]{PAL rule sheet} carry the distinct cited
whole-language reduction result; this calculation does not replace its
adequacy boundary \cite[p.~189]{VanBenthemLiuSmets2021Logico}.
The related muddy-children result in \cref{thm:p08-muddy-children}
uses repeated publicly observable answers under a fixed truthful,
synchronous protocol. Its assumptions explain why announcements of
ignorance can eliminate possibilities. Distributed knowledge, as in
\cref{def:p08-distributed-knowledge}, pools information by intersecting
accessibility relations; common knowledge instead quantifies over
arbitrarily long paths. Neither pooled information nor a finite nesting
of ``everyone knows'' can silently replace the common-knowledge
condition in a protocol.

\begin{hbdefinition}{def:p08-rat26-soft-update}{Finite plausibility and two upgrades}
For one agent let \(W\) be finite and nonempty, with a total preorder
\(\preceq\), where earlier worlds are more plausible. Write
\(\min P\) for the best worlds in a nonempty \(P\subseteq W\).
Knowledge \(K\phi\) requires \(\phi\) at every world; belief
\(B\phi\) requires it only at \(\min W\). Conditional belief
\(B^\chi\psi\) requires \(\psi\) at \(\min\llbracket\chi\rrbracket\)
when that set is nonempty, and is vacuously true otherwise.
For factual information with truth set \(P\ne\varnothing\), the
conservative upgrade \(\uparrow\phi\) promotes just \(\min P\)
to a new first tier, leaving the ordering of all other worlds unchanged.
The radical upgrade \(\Uparrow\phi\) places every \(P\)-world
before every non-\(P\)-world and preserves the order inside each class.
Neither deletes worlds or changes atomic valuations. For empty \(P\)
we stipulate the identity operation.
\end{hbdefinition}

These are bounded finite versions of the operations in
\cite[pp.~257--261]{BaltagSmets2011TruthUpgrades}; they make explicit
the chapter's contrast between hard and soft information. Fallibility
is represented by leaving alternatives available, not by pretending
that an unreliable report is known to be true.

\begin{hbtheorem}{thm:p08-rat26-soft-update}{proved=annex:p08.rat2.6.soft-update}{Agreement on current belief does not identify an update}
For factual \(\phi\), \(\neg K\neg\phi\) implies
\([\uparrow\phi]B\phi\) and \([\Uparrow\phi]B\phi\).
Knowledge of factual propositions is unchanged by either operation.
Nevertheless they can yield different conditional beliefs. Take
\(a\prec b\prec c\), \(P=\{b,c\}\),
\(\llbracket q\rrbracket=\{c\}\), and
\(\llbracket\chi\rrbracket=\{a,c\}\). After conservative upgrade
the order is \(b\prec a\prec c\), with \(\neg B^\chi q\);
after radical upgrade it is \(b\prec c\prec a\), with \(B^\chi q\).
Both have \(B\phi\) but not \(K\phi\).
\end{hbtheorem}

Repeated learning raises a further question: under which truthfulness,
fairness and information-supply conditions do beliefs approach the
truth? A single-upgrade calculation does not prove a convergence theorem
for an infinite stream. The source's belief-change programme also
considers evidence, inquisitive issues, changing preferences and
commands affecting obligations. Those enrich the state and the update
type; they cannot all be modeled by silently reusing truthful factual
restriction \cite[pp.~189--190]{VanBenthemLiuSmets2021Logico}.

Classical consequence itself is monotonic: if \(\Gamma\models\psi\),
then \(\Gamma\cup\Delta\models\psi\). A default such as ``birds
normally fly'' can instead be defeated by learning that this bird is a
penguin. Belief revision changes a represented state and may withdraw
accepted conclusions while retaining classical inference within that
state. A nonmonotonic consequence relation changes the inferential
relation instead. Nonclassical and substructural logics offer still
other choices. These are distinct mechanisms, not different names for
the same public-announcement operation.

\section{Interactive agency and representation-sensitive difficulty}

The source's extensive game has agent A choose between stopping with
payoff \((1,0)\) and allowing E to choose between \((0,100)\) and
\((99,99)\). With those payoffs, perfect information and sequential
maximization of each player's own payoff, E chooses 100 over 99; A
therefore chooses 1 over 0 and the outcome is \((1,0)\). Both players
would prefer \((99,99)\) to that outcome. The example displays a
tension between individual continuation incentives and a jointly better
outcome; it is not a theorem that all rational agents accept the stated
payoffs, expectations or solution concept
\cite[pp.~189--190]{VanBenthemLiuSmets2021Logico}.

Information about another agent's information affects strategy. Higher-
order knowledge in a model is not, however, empirical evidence that
human players perform that recursion. Groups, coalitions and distributed
architectures can supply capabilities absent in a lone agent; they can
also entrench mistakes. Elimination of dominated strategies, regret
minimization and Boolean-game representations are source-cited methods
within this broader interactive programme, not interchangeable universal
tests of rationality \cite[pp.~190--191]{VanBenthemLiuSmets2021Logico}.

\begin{hbdefinition}{def:p08-rat26-sabotage}{Reachability sabotage with explicit edges}
An input is a finite directed or undirected multigraph, with parallel
edges explicitly listed, an initial vertex and a target. Runner moves
first along one available edge. Unless the target has been reached,
Blocker then deletes one available edge anywhere in the graph. Turns
alternate with perfect information and no passing. Runner wins upon
reaching the target, including at the initial position; otherwise Runner
loses when no outgoing edge remains. In the undirected case deleting an
edge removes that same connection in both directions. The input length
counts the listed edges, not a succinct binary encoding of arbitrarily
large multiplicities.
\end{hbdefinition}

\begin{hbtheorem}{thm:p08-rat26-sabotage-space}{proved=annex:p08.rat2.6.sabotage-space}{A polynomial-space search for sabotage reachability}
The winning-strategy problem in \cref{def:p08-rat26-sabotage} is
decidable in polynomial space in the explicit input size.
\end{hbtheorem}

\begin{hbtheorem}{thm:p08-rat26-sabotage-complete}{cited=Rohde2006DynamicStructures}{The interactive complexity boundary}
For the explicit finite multigraph reachability game just specified,
deciding whether Runner has a winning strategy is PSPACE-complete,
including instances with a single target; the result applies to
undirected and directed arenas.
\end{hbtheorem}

Rohde's Definitions 2.1--2.2 fix the size and move convention; Theorem
2.20 and Remark 2.21 supply the cited lower-bound result
\cite[pp.~24--26, 43--44]{Rohde2006DynamicStructures}. The external
hardness proof reduces quantified Boolean formulas, arranging assignment
choices and verification while controlling Blocker's opportunity to
delete edges outside the intended local passage. A naive assignment
chain is insufficient precisely because deletion is global. The local
annex proves only the polynomial-space upper bound, not that reduction.
Thus the chapter's comparison with ordinary static reachability has an
identified formal boundary: a path that exists initially need not be
enforceable against a changing arena. Static reachability is polynomial-
time decidable; the interactive problem has the stated completeness
classification. An unconditional separation of P from PSPACE is not
claimed, nor is this classification transferred to every sabotage
variant or compressed graph representation.

\section{High-level reasoning and low-level population dynamics}

The chapter challenges a restriction of rationality to reflective,
linguistically articulated deliberation. A population of fixed behavioral
types can have the same aggregate frequencies as a deliberating agent's
randomized choices, although the mechanisms differ. Its Hawk--Dove
illustration needs a precise payoff convention before drawing a stable-
mixture conclusion \cite[p.~191]{VanBenthemLiuSmets2021Logico}.

\begin{hbdefinition}{def:p08-rat26-hawk-dove}{A specified Hawk--Dove invasion test}
Let resource value \(V>0\) and injury cost \(C>0\). The row
player's payoff matrix is
\[
\begin{array}{c|cc}
 &H&D\\\hline
H&(V-C)/2&V\\
D&0&V/2
\end{array}
\]
with random independent pairing. Write \(u(x,z)\) for the expected
payoff when the row player uses Hawk with probability \(x\), against
Hawk frequency \(z\). A resident \(x\in[0,1]\) is evolutionarily
stable here if for each mutant \(y\ne x\) there is \(\varepsilon_0>0\)
such that, for \(0<\varepsilon<\varepsilon_0\),
\[
 u(x,(1-\varepsilon)x+\varepsilon y)
 >u(y,(1-\varepsilon)x+\varepsilon y).
\]
This is a defined invasion criterion, not a claim that every population
trajectory converges, or that fitness establishes a moral obligation.
\end{hbdefinition}

\begin{hbtheorem}{thm:p08-rat26-hawk-dove}{proved=annex:p08.rat2.6.hawk-dove}{The cost condition for the stable mixture}
For the specified matrix, the unique evolutionarily stable resident
is \(x^*=V/C\) when \(0<V<C\), and pure Hawk \(x=1\) when
\(V\ge C\), including equality. In particular an interior stable
mixture requires injury cost to exceed resource value.
\end{hbtheorem}

This corrects the direction of the chapter's unqualified inequality for
the usual matrix. The distinction between a mixed strategy and a
population mixture remains valuable: identical expected payoffs do not
show identical reasoning. The source's analogy between law-abiding
agents and free riders similarly poses a social-dynamic question; it
does not derive which laws ought to be obeyed. Ordinary agents can
combine habits, learning and occasional deliberation, rather than
belonging permanently to one ``high'' or ``low'' category.

\begin{hbdefinition}{def:p08-rat26-influence}{The source's synchronous influence graph}
On vertices \(\{1,2,3,4\}\), the outgoing consultation neighborhoods
are \(N(1)=\{2,3\}\), \(N(2)=\{4\}\), \(N(3)=\{2\}\),
\(N(4)=\{3\}\). An arrow \(s\to t\) means that \(s\) consults
\(t\). Given active set \(X\), all agents update simultaneously by
\[
 T(X)=\{s:N(s)\subseteq X\}.
\]
Thus the direction is not silently reversed into an incoming-neighbor
convention. On a general finite graph, an empty consultation set satisfies
this condition vacuously unless a different rule is declared.
\end{hbdefinition}

\begin{hbtheorem}{thm:p08-rat26-network}{proved=annex:p08.rat2.6.network}{Cycles, transients and the structure being compared}
The four source examples evolve as follows, with the indicated cycle
repeated thereafter:
\[
\begin{array}{c|l}
X_0&X_1,X_2,\ldots\\\hline
\{1\}&\varnothing,\varnothing,\ldots\\
\{2\}&\{3\},\{4\},\{2\},\ldots\\
\{1,2\}&\{3\},\{4\},\{2\},\ldots\\
\{1,2,3\}&\{1,3,4\},\{2,4\},\{2,3\},\{1,3,4\},\ldots
\end{array}
\]
Every deterministic update on a finite graph is eventually periodic,
but need not stabilize. The full directed graph above has only the
identity automorphism. Consequently its different singleton-marked
states are not isomorphic as full marked-network models, although the
induced active singleton subgraphs are all isomorphic.
\end{hbtheorem}

The source Figure 2.6.2 illustrates a deterministic dynamics, not a
measured social contagion law \cite[pp.~191--192]{VanBenthemLiuSmets2021Logico}.
Even without any agent proving a theorem, a logical or mathematical
analysis can identify the update, its invariants and its cycles. This
is the chapter's important foundational role for logic at a low level.
Finiteness is essential to the repetition argument; an infinite state
space can have a nonrepeating trajectory. A convergence claim would
require an additional property, not merely synchronous updating.

The source then juxtaposes such rules with epistemic models of information
cascades. In the reviewed cascade setting, higher-order knowledge alone
does not stop the cascade, while adequate communication of evidence can.
That is a claim within a protocol and informational model, not a theorem
that no form of higher-order reasoning or intervention can ever prevent
a cascade. No population frequencies or causal effect sizes are supplied
by this chapter. Network structure, evidence access and strategic
expectations are separate explanatory ingredients.

\section{Machine learning, probability and the remaining objections}

The neural-network discussion describes adjustable weights, thresholded
processing and optimization against labeled examples. Its challenge to
logic is representational: the learned mechanism may have no immediately
recognizable sentence-level logical language or obvious inference
operations. It is not an impossibility proof for logical analysis.
Logic can analyze a learning theory at a foundational level, connect to
causal explanations, or investigate extracted representations, even when
the device itself is not transparently executing a deductive calculus.
The source argues for cooperation between approaches, with hybrid and
explainable systems as research directions, not a demonstrated universal
advantage over all alternatives
\cite[pp.~192--193]{VanBenthemLiuSmets2021Logico}.

\begin{hbtheorem}{thm:p08-rat26-learning-limits}{proved=annex:p08.rat2.6.learning-limits}{Two limits of an optimization description}
Gradient descent alone does not guarantee convergence: on \(J(t)=t^2\),
the fixed-step iteration with \(t_0\ne0\) diverges in magnitude when
the step size \(\eta>1\), and oscillates when \(\eta=1\).
Zero training error alone does not determine held-out predictions:
the functions \(f_c(x)=cx\) on \(\{0,1\}\), trained only on the
labeled example \((0,0)\), all have zero squared training error but
have squared error \(c^2\) at the held-out labeled example \((1,0)\).
\end{hbtheorem}

These countercases qualify the chapter's informal convergence and
generalization language. They do not deny successful learning under
suitable loss, data, architecture and optimization assumptions. Nor do
they substitute for an empirical comparison of trained systems. A
description of a cost-minimizing method must specify what is actually
guaranteed, for which data and procedure. P12's result and encoding audits
apply here as methodological questions, but a trained predictor is not
thereby the same artifact as a formal proof certificate.

Probability supplies another complementary language. The source
distinguishes qualitative epistemic possibilities from quantitative
degrees of belief, and recalls the historical connections of logic with
probability through Boole and de Finetti. Probability can represent
uncertainty about the world or information; a mixed strategy instead
randomizes action. Its probabilities need not be the agent's credences
about a proposition. Probabilistic parsing can combine remembered data
with rule-like structure; qualitative belief and probability-comparison
languages can also be related without becoming identical. Neither the
presence of numbers nor their absence establishes that an approach is
the only theory of rational agency
\cite[pp.~193--194]{VanBenthemLiuSmets2021Logico}.

The conclusion retains three open boundaries. First, idealized models
of logically omniscient knowledge can be valuable for information
structure without being cognitive process models. Second, abstract
programmability does not establish a feasible implementation or remove
resource constraints. Third, formal, normative and descriptive uses
of a rationality model may pull apart. The chapter expressly leaves
empirical cognition outside its main scope; the semantic proofs here
must not be relabeled as psychological evidence. A universal canonical
model of rational agency is posed as a question, not supplied, and its
absence is not proved inevitable. The source's contribution is a
plural but connected investigation of inference, information and
interaction, with explicit operations and limits rather than a single
norm imposed on all agents.
